\documentclass[twocolumn]{article}
\usepackage[T1]{fontenc}
\usepackage[utf8]{inputenc}
\usepackage{amsmath,amssymb}
\usepackage{graphicx}
\usepackage{booktabs}
\usepackage{hyperref}
\usepackage{geometry}
\usepackage{array}
\usepackage{caption}
\usepackage{subcaption}
\usepackage{xcolor}
\usepackage{microtype}
\geometry{margin=1in}

\title{\textbf{A Hybrid Secure Big-Data Pipeline Combining\\
Variational Autoencoder Compression, AES-256-GCM\\
Authenticated Encryption, and Post-Quantum\\
ML-KEM/Kyber-1024 Key Encapsulation}}

\author{[Author 1]$^1$, [Author 2]$^2$, [Author 3]$^1$\\[6pt]
$^1$Laboratory of Computer Science, Data Science and Artificial Intelligence,\\
National Higher Polytechnic School of Douala, University of Douala,\\
P.O.\ Box 8580, Douala, Cameroon\\[2pt]
$^2$Laboratory of Electrical, Electronics, Automation and Telecommunication,\\
National Higher Polytechnic School of Douala, University of Douala,\\
P.O.\ Box 8580, Douala, Cameroon\\[4pt]
\textit{Correspondence:} [Author 1], [email]}

\date{Submitted to \textit{Complex \& Intelligent Systems} --- [Submission Date]}

\begin{document}
\maketitle

\noindent\textbf{Keywords:} Data security pipeline; Variational autoencoder; AES-256-GCM;
Post-quantum cryptography; ML-KEM/Kyber; Authenticated encryption;
Heterogeneous data; Reproducible research

% ─────────────────────────────────────────────────────────────
\begin{abstract}
Securing heterogeneous data requires simultaneously reducing data volumes,
guaranteeing confidentiality and integrity, and maintaining acceptable performance
as data scale grows. We propose and evaluate a reproducible hybrid pipeline that
sequentially combines a Variational Autoencoder (VAE) for pre-encryption
dimensionality reduction, AES-256-GCM authenticated encryption, and a
post-quantum ML-KEM/Kyber-1024 key encapsulation mechanism. The pipeline is
evaluated on a real-world dataset of 126 files totalling 3.79~GB (75~JPEG images
and 50~binary MP4 files), using micro-benchmarks, NIST CAVP test-vector
validation (775/775 vectors passed, 100\%), a scalability study over 7 data
points from 0.60~GB to 3.79~GB, and a full round-trip fidelity analysis.

A central design decision acknowledged throughout this work is that VAE
compression is inherently lossy: the pipeline protects the latent representation
$\hat{z}$ of the original data $D$, not $D$ itself. The recovered output after
full round-trip decryption is an approximation $\hat{D} \approx D$. This makes
the pipeline appropriate for scenarios where approximate reconstruction is
acceptable --- such as exploratory data analysis, similarity search, and preview
streaming --- and explicitly unsuitable for archival, financial, or diagnostic
applications requiring bit-exact reconstruction.

Results show that (i)~the direct cryptographic cost (AES-GCM + ML-KEM combined)
accounts for 97.6--100.2\% of measured per-file processing latency (median
34.24~ms for all file types combined), confirming that VAE inference and I/O are
the dominant unaccounted costs in a production deployment; (ii)~throughput across
four ablation configurations ranges from 23.9~MB/s (gzip+AES+KEM) to
993.6~MB/s (AES-only), with ML-KEM overhead (12.6--19.9~ms per operation)
remaining modest relative to I/O at larger file sizes; (iii)~scalability is
essentially flat over the measured range (slope $\beta = -0.09$~MB/s/GB,
$R^2 < 0.01$), indicating no throughput collapse in the medium-scale regime; and
(iv)~an ablation study confirms that each pipeline layer contributes meaningfully
to the overall security-performance profile. The primary contribution is a
system-level integration and reproducible joint evaluation of three complementary
security layers --- learned compression, standardised authenticated encryption,
and post-quantum key encapsulation --- on real multi-format data.
\end{abstract}

% ─────────────────────────────────────────────────────────────
\section{Introduction}

The digitalisation of industrial, medical, and governmental systems creates
complex, heterogeneous, and rapidly growing data ecosystems whose security remains
an open engineering challenge. Data generated by these systems spans images,
binary executables, relational tables, and unstructured logs, each with distinct
statistical properties that affect both compression efficiency and cryptographic
cost.

Two independent technological pressures are reshaping the design of data security
systems. First, standardised symmetric cryptography (AES-256) and hash-based
authentication (HMAC-SHA256) are computationally intensive when applied to large,
uncompressed data volumes. Reducing data volume prior to encryption lowers the
computational footprint of cryptographic operations without sacrificing security
guarantees. Second, advances in quantum computing threaten the long-term security
of currently deployed asymmetric schemes. Several independent analyses establish
that Shor's algorithm, running on a cryptographically relevant quantum computer
(CRQC) --- a fault-tolerant quantum device capable of breaking 2048-bit RSA
within hours --- would compromise both RSA and elliptic-curve cryptography, making
migration to post-quantum primitives a long-term engineering necessity
\cite{natacad2019, bernstein2017}.

Generative neural compression, specifically variational autoencoders (VAEs),
offer a complementary lever: reducing data volume prior to encryption lowers the
computational footprint of cryptographic operations on heterogeneous pipelines
without assigning the AI component any cryptographic role. The VAE operates
strictly as a pre-processor; all security properties are provided by the
downstream cryptographic layer.

This paper makes the following contributions:
\begin{enumerate}
  \item \textbf{A modular three-layer security architecture} combining VAE
  compression, AES-256-GCM authenticated encryption, and ML-KEM/Kyber-1024
  post-quantum key encapsulation, with each layer independently auditable and
  replaceable.
  \item \textbf{Explicit fidelity transparency}: the system protects the latent
  representation $\hat{z}$, not the original plaintext $D$. Round-trip fidelity
  is quantified explicitly.
  \item \textbf{Reproducible benchmarking} on 3.79~GB of real heterogeneous data
  (75~JPEG images, 50~binary MP4 files, 1 JSON text file; 126 files total),
  including per-layer latency breakdown, throughput scalability, and NIST CAVP
  correctness validation.
  \item \textbf{Ablation study} isolating the contribution of each pipeline layer
  and providing baselines for future comparison.
  \item \textbf{Informal security composition argument} establishing that, under
  standard hardness assumptions, the pipeline inherits IND-CCA2 security from its
  cryptographic components, with the VAE contributing no cryptographic guarantee.
\end{enumerate}

The remainder of the paper is structured as follows. Section~2 provides
background on heterogeneous data and cryptographic foundations. Section~3 surveys
related work and identifies the research gap. Section~4 presents the pipeline
architecture, threat model, formal metric definitions, and security composition.
Section~5 reports experimental results. Section~6 discusses findings in context.
Section~7 concludes and outlines future directions.

% ─────────────────────────────────────────────────────────────
\section{Background}

\subsection{Heterogeneous Big Data}

Massive heterogeneous data, broadly referred to as ``big data'' in the
literature, is defined not by an absolute volumetric threshold but by the
inadequacy of traditional management tools given its intrinsic properties:
heterogeneity of formats, generation velocity, and scale of processing
requirements. The three foundational dimensions --- Volume, Variety, and Velocity
--- were first introduced by Laney \cite{laney2001} and remain the dominant
characterisation framework. Subsequent work has extended this taxonomy to include
Veracity and Value \cite{gandomi2015}. For the purposes of this paper, the
Variety dimension is most significant: the pipeline must process images and binary
files without format-specific preprocessing, which creates a direct challenge for
any compression scheme that assumes a particular data distribution.

A critical implication of high-entropy data --- common among binary executables,
video files, and encrypted archives --- is that classical lossless compressors
(gzip, zstd, bzip2) achieve minimal compression (CR $\approx$ 95--100\% of
original size). This motivates the use of a learned compressor (VAE), which can
project high-dimensional inputs into a low-dimensional latent space regardless of
entropy, at the cost of introducing reconstruction error.

\subsection{AES-256-GCM Authenticated Encryption}

The Advanced Encryption Standard (AES), standardised by NIST in 2001, operates
on 128-bit blocks using key sizes of 128, 192, or 256 bits. Its $2^{256}$ key
space provides a security level of 256 bits against classical brute-force attacks;
Grover's quantum search algorithm reduces this effective security to $2^{128}$
bits, which remains computationally intractable at foreseeable quantum computing
scales \cite{grover1996}.

Galois/Counter Mode (GCM) extends AES with a built-in Galois field authentication
tag, producing Authenticated Encryption with Associated Data (AEAD): a
single-pass operation that simultaneously guarantees confidentiality and integrity.
AES in GCM mode is natively parallelisable, as each counter block is encrypted
independently. NIST SP 800-38D \cite{dworkin2007} specifies the full construction.
A critical implementation requirement is nonce uniqueness: reuse of a 96-bit
nonce under the same key catastrophically breaks both confidentiality and
integrity \cite{aviram2016}.

\subsection{Post-Quantum Key Encapsulation: ML-KEM/Kyber-1024}

CRYSTALS-Kyber, standardised by NIST in August 2024 as Module Lattice Key
Encapsulation Mechanism (ML-KEM, FIPS~203) \cite{fips203}, is the primary
NIST-selected post-quantum key encapsulation mechanism. Its security rests on the
Module Learning With Errors (MLWE) hardness assumption: given a matrix $A \in
\mathbb{Z}_q^{k \times k}$, vector $t \in \mathbb{Z}_q^k$ (derived from secret
$s$ and error $e$), it is computationally infeasible to recover $s$ from the pair
$(A,\,t = As + e \bmod q)$.

The three security parameter sets --- Kyber-512 (AES-128 equivalent), Kyber-768
(AES-192 equivalent), and Kyber-1024 (AES-256 equivalent) --- differ in module
rank $k$ (2, 3, and 4 respectively), with corresponding increases in public-key
size, ciphertext size, and encapsulation latency. This paper uses Kyber-1024 to
match the AES-256-GCM security level, providing consistent 256-bit post-quantum
security throughout the pipeline.

% ─────────────────────────────────────────────────────────────
\section{Related Work and Research Gap}

\begin{table}[t]
\caption{Comparison of related pipeline approaches.}
\label{tab:related}
\centering\small
\begin{tabular}{lllll}
\toprule
Reference & Compression & Encryption & PQ & Data \\
\midrule
Zhang et al.~\cite{zhang} & Classical (gzip) & AES-128-CBC & \texttimes & Images \\
Fu et al.~\cite{fu} & None & AES-256-GCM & \texttimes & IoT logs \\
Nugroho et al.~\cite{nugroho} & DCT & AES-128-GCM & \texttimes & Images \\
Thenmozhi et al.~\cite{thenmozhi} & CNN & ChaCha20 & \texttimes & Medical \\
Alsubai et al.~\cite{alsubai} & Autoencoder & RSA & \texttimes & Mixed \\
Wang \& Lo~\cite{wanglo} & Joint AE & Embedded & \texttimes & Images \\
\midrule
\textbf{This work} & VAE & AES-256-GCM & ML-KEM-1024 & 3.79~GB hetero. \\
\bottomrule
\end{tabular}
\end{table}

Wang and Lo \cite{wanglo} propose a joint autoencoder-encryption scheme in which
encryption is embedded directly within the latent representation, achieving
end-to-end co-optimisation at the cost of non-standard, non-auditable
cryptographic primitives. Thenmozhi et al.\ \cite{thenmozhi} apply convolutional
autoencoders to medical image compression before ChaCha20 encryption, showing
compression gains but no formal security composition. Zhang et al.\ \cite{zhang}
and Fu et al.\ \cite{fu} use classical (non-neural) compression, providing strong
throughput baselines but no post-quantum key management. Fall et al.\ \cite{fall}
systematise hybrid classical/post-quantum strategies, demonstrating that hybrid
KEM-DEM constructions maintain IND-CCA2 security when the KEM and DEM are
independently IND-CCA2 secure. Alsubai et al.\ \cite{alsubai} combine autoencoder
compression with RSA key encapsulation but do not address the post-quantum threat.
Nugroho et al.\ \cite{nugroho} demonstrate that applying compression before
encryption reduces ciphertext size on image workloads but do not evaluate binary
data, and do not quantify per-layer latency contributions.

Three gaps motivate this work: (i) no existing work combines all three layers
(neural compression + AEAD + post-quantum KEM) using NIST-standardised
primitives; (ii) existing pipelines report aggregate throughput without isolating
each layer's cost; (iii) no existing work provides a complete, open-source,
dataset-available pipeline that a practitioner can replicate on their own
infrastructure.

% ─────────────────────────────────────────────────────────────
\section{Pipeline Architecture}

\subsection{Three-Layer Sequential Design}

The pipeline applies three sequential transformations to input data $D$:

\medskip
$D \;\longrightarrow\; [\text{VAE Encoder}] \;\longrightarrow\; z \;\longrightarrow\; [\text{AES-256-GCM}] \;\longrightarrow\; c \;\longrightarrow\; [\text{ML-KEM Encaps}] \;\longrightarrow\; (c,\,K_\text{enc})$
\medskip

\textbf{Layer 1 (VAE Compression).} The encoder $f_\theta$ maps $D$ to a latent
vector $z = f_\theta(D) \in \mathbb{R}^k$. For 32$\times$32 RGB image inputs,
$k = 128$ (latent dim). The latent vector is serialised to float32 bytes (512
bytes), yielding the compressed representation $S(D)$.

\textbf{Layer 2 (AES-256-GCM).} A session key $K \in \{0,1\}^{256}$ is sampled
uniformly at random. The ciphertext is $c = \text{AES-GCM}_K(S(D),\,N,\,\text{AAD})$,
where $N$ is a 96-bit nonce and AAD contains object metadata. AES-GCM appends a
128-bit authentication tag $\tau$, guaranteeing that any modification to $c$ is
detected with probability $1 - 2^{-128}$.

\textbf{Layer 3 (ML-KEM-1024).} The recipient's public key $(pk, sk)$ is
generated via ML-KEM.KeyGen(). The session key $K$ is encapsulated as
$(K,\,K_\text{enc}) = \text{ML-KEM.Encaps}(pk)$. The transmitted bundle is
$(c,\,\tau,\,K_\text{enc},\,N,\,\text{AAD})$.

\textbf{Decapsulation.} The recipient recovers $K$ via
ML-KEM.Decaps$(sk,\,K_\text{enc})$, decrypts $c$ to obtain $S(D)$ (verified by
$\tau$), deserialises $z$, and decodes $\hat{D} = g_\phi(z)$ via the VAE decoder.

\smallskip
\noindent\textbf{Critical limitation:} The AES-GCM authentication tag $\tau$
guarantees the integrity of $S(D)$, not of $D$. Because the VAE is lossy,
$\hat{D} \neq D$ in general. The pipeline is appropriate only for use cases where
approximate reconstruction is acceptable.

\subsection{ML-KEM-1024 Algorithm Specification}

All notation follows FIPS~203 \cite{fips203}. Parameters: module rank $k = 4$,
modulus $q = 3329$, polynomial degree $n = 256$, noise parameters $\eta_1 = 2$,
$\eta_2 = 2$, compression parameters $(d_u, d_v) = (11, 5)$.

\textbf{KeyGen:} (i)~Sample $A \leftarrow \mathbb{Z}_q^{k\times k}$ from XOF(seed~$\rho$);
(ii)~sample secret $s$, error $e$ from centred binomial distribution $B_{\eta_1}$;
(iii)~compute $t = As + e \bmod q$; (iv)~public key: $pk = (A,t)$, private key: $sk = s$.

\textbf{Encaps:} (i)~Sample $m \leftarrow \{0,1\}^{256}$ uniformly; (ii)~derive
$(\bar{K}, r) = G(m \,\|\, H(pk))$; (iii)~sample $r_1,r_2,e_2$ from
$B_{\eta_1},B_{\eta_2}$; (iv)~compute $u = A^T r_1 + r_2 \bmod q$, $v = t^T r_1
+ e_2 + \lfloor q/2 \rfloor \cdot m \bmod q$; (v)~return $(K,\,K_\text{enc} =
\text{Compress}(u,v))$.

\textbf{Decaps:} (i)~Decompress $(u,v)$ from $K_\text{enc}$; (ii)~recover
$m' = \text{Decompress}(v - s^T u \bmod q)$; (iii)~re-encapsulate with $m'$ and
verify implicitly (implicit rejection on mismatch); (iv)~return $K = H(m' \,\|\,
H(pk))$.

Kyber-1024 targets NIST security category 5 ($\geq 2^{256}$ classical operations,
$\geq 2^{128}$ quantum operations for the best known MLWE attack \cite{fips203}).

\subsection{Threat Model}

\begin{table}[t]
\caption{Threat model summary.}
\label{tab:threat}
\centering\small
\begin{tabular}{p{1.4cm}p{2.0cm}p{2.2cm}p{1.8cm}}
\toprule
Threat & Adversary & Countermeasure & Guarantee \\
\midrule
T1: Eavesdrop. & Passive observer & AES-256-GCM + ML-KEM-1024 & 256-bit classical / 128-bit quantum \\
T2: Tampering & Active MITM & AES-GCM $\tau$ & $2^{-128}$ forgery prob. \\
T3: HNDL & Quantum future & ML-KEM-1024 (MLWE) & NIST Category 5 \\
T4: Latent infer. & Observer with $z$ & $z$ encrypted under AES & Reduces to T1 \\
\bottomrule
\end{tabular}
\end{table}

\textbf{Out of scope:} Key management infrastructure, certificate authorities,
side-channel attacks on hardware implementations, and adversarial perturbations
to the VAE decoder.

\subsection{Security Composition}

Under the IND-CCA2 security of ML-KEM-1024 (Assumption H1: MLWE hardness), the
IND-CPA security of AES-256 (Assumption H2: AES pseudorandomness), and the
unforgeability of GCM-GHASH (Assumption H3: collision resistance of the Galois
field multiplication), the hybrid KEM-DEM construction inherits IND-CCA2
security. This follows from the hybrid KEM-DEM composition theorem of Shoup
\cite{shoup2001}, confirmed by the Cramer-Shoup framework \cite{cramer2003}.

\textbf{Scope limitation:} This composition applies strictly to the cryptographic
components. The VAE provides no cryptographic guarantee. The pipeline therefore
provides IND-CCA2 protection of the latent representation $\hat{z}$, not of the
original data $D$.

\subsection{Evaluation Metrics}

\begin{itemize}
  \item \textbf{Compression ratio (CR)}: $|S(D)|/|D|$ (lower is better).
  \item \textbf{Reconstruction fidelity}: SSIM$(\hat{D}, D)$ and PSNR$(\hat{D}, D)$ for images; byte-level Hamming distance for non-image data.
  \item \textbf{Per-layer latency}: $L_\text{total} = L_\text{VAE} + L_\text{AES} + L_\text{KEM} + L_\text{IO}$, reported as median [IQR].
  \item \textbf{Throughput}: $\Theta = |D|/L_\text{total}$ (MB/s). Scalability regression: $\Theta(s) = \alpha + \beta s$.
  \item \textbf{NIST CAVP pass rate}: fraction of NIST-provided test vectors producing bit-exact outputs.
\end{itemize}

% ─────────────────────────────────────────────────────────────
\section{Experimental Results}

\subsection{Dataset and Experimental Setup}

\textbf{Dataset D1} comprises 126 real files totalling 3.79~GB, used for all
micro-benchmarks, latency analysis, and scalability evaluation. The composition
is detailed in Table~\ref{tab:d2}. Note that MP4 video files constitute 87.1\%
of the dataset by volume; classical lossless compressors achieve minimal
compression on this high-entropy content.

\begin{table}[t]
\caption{D1 dataset composition by file type.}
\label{tab:d2}
\centering\small
\begin{tabular}{lrrrr}
\toprule
File type & Count & Total (GB) & Mean (MB) & Fraction (\%) \\
\midrule
Binary (MP4) & 50 & 3.30 & 67.3 & 87.1 \\
Images (JPEG) & 75 & 0.49 & 6.6 & 12.9 \\
Text (JSON) & 1 & $<$0.01 & 0.1 & $<$0.1 \\
\midrule
\textbf{Total D1} & 126 & 3.79 & 30.9 & 100.0 \\
\bottomrule
\end{tabular}
\end{table}

\begin{table}[t]
\caption{Hardware and software configuration.}
\label{tab:config}
\centering\small
\begin{tabular}{ll}
\toprule
Component & Specification \\
\midrule
Platform & Windows 11 Pro (local workstation) \\
Python & 3.13.7 \\
PyTorch & 2.12+ (CPU-only) \\
cryptography & hazmat AES-GCM (FIPS 197) \\
ML-KEM library & kyber-py shim (liboqs interface) \\
Kyber variant & ML-KEM-1024 (Kyber-1024) \\
\bottomrule
\end{tabular}
\end{table}

\textbf{Implementation note:} The pipeline runs on a CPU-only workstation under
Windows~11. The ML-KEM implementation uses a pure-Python \texttt{kyber-py} shim
that is functionally identical to \texttt{liboqs} but slower than native
C-compiled liboqs. Absolute latency figures should be treated as indicative of
relative cost proportions rather than hardware-optimised absolute performance.
A Linux/native-liboqs deployment (the production target) will achieve
substantially lower ML-KEM latency.

\subsection{NIST CAVP Correctness Validation}

\begin{table}[t]
\caption{NIST CAVP validation results. All pass rates are 100\%.}
\label{tab:cavp}
\centering\small
\begin{tabular}{lrrrr}
\toprule
Test suite & Total & Pass & Fail & Pass rate \\
\midrule
AES-256-GCM Encrypt KAT & 375 & 375 & 0 & 100.0\% \\
AES-256-GCM Decrypt (incl.\ FAIL) & 300 & 300 & 0 & 100.0\% \\
ML-KEM-1024 Keygen/Encaps/Decaps & 100 & 100 & 0 & 100.0\% \\
\midrule
\textbf{Total} & 775 & 775 & 0 & 100.0\% \\
\bottomrule
\end{tabular}
\end{table}

All 775 NIST CAVP test vectors for AES-256-GCM and ML-KEM-1024 passed
(Table~\ref{tab:cavp}), confirming that the cryptographic implementations are
correct and interoperable with reference implementations.

\subsection{AES-256-GCM Micro-Benchmark}

AES-256-GCM throughput was measured across plaintext sizes of 1~KB to 64~KB:

\begin{itemize}
  \item 1~KB: \textbf{109~MB/s} (encrypt), 112~MB/s (decrypt)
  \item 4~KB: \textbf{532~MB/s} / 537~MB/s
  \item 16~KB: \textbf{997~MB/s} / 1,027~MB/s
  \item 64~KB: \textbf{1,387~MB/s} / 1,344~MB/s
\end{itemize}

Throughput increases sharply as payload size grows, reflecting the amortisation
of per-operation Python overhead across larger payloads. At 64~KB, AES-GCM
exceeds 1~GB/s in both directions.

\subsection{ML-KEM Micro-Benchmark}

\begin{table}[t]
\caption{ML-KEM key encapsulation latency by security level.}
\label{tab:kem}
\centering\small
\begin{tabular}{lrrr}
\toprule
Algorithm & KeyGen ($\mu$s) & Encaps ($\mu$s) & Decaps ($\mu$s) \\
\midrule
Kyber-512 & 3,827 & 5,297 & 7,661 \\
Kyber-768 & 7,965 & 9,927 & 13,914 \\
\textbf{Kyber-1024} & \textbf{12,590} & \textbf{14,924} & \textbf{19,897} \\
ML-KEM-512 & 4,709 & 6,581 & 9,481 \\
\bottomrule
\end{tabular}
\end{table}

The selected Kyber-1024 variant requires a total key-exchange cycle (KeyGen +
Encaps + Decaps) of approximately 47.4~ms under our pure-Python implementation.
This overhead is per-message, not per-byte, meaning it becomes negligible for
large files but significant for small files (e.g.\ a 7~KB JPEG image where
47~ms of KEM cost exceeds 7~ms of AES cost).

\subsection{VAE Compression Quality}

The VAE was trained for 2 epochs on 75 real JPEG images (fast validation run;
production deployment targets 20 epochs). After 2 epochs:
\begin{itemize}
  \item Validation MSE loss: 0.036
  \item Validation PSNR: 14.48~dB
\end{itemize}

The low PSNR (14.48~dB vs.\ the SSIM $\geq 0.90$ threshold) reflects the
prototype status of the current model trained with minimal epochs. The VAE
encodes a 32$\times$32$\times$3 input (3,072 bytes) to a 128-dimensional latent
vector (512 bytes float32), achieving a compression factor of 6$\times$ on the
resized input. Applied to the original multi-megabyte source files (which are
first downsampled to 32$\times$32), the effective compression is extreme and
extremely lossy; this is appropriate only for thumbnail-quality preview and
similarity-search use cases.

\textbf{Classical compressors on image data:} Lossless compressors (gzip, bz2,
lz4, brotli, zstd) applied to already-compressed JPEG files achieve a compression
ratio of approximately 1.0$\times$ (SSIM~=~1.00, PSNR~=~99.0~dB), confirming
that JPEG files are already near-incompressible by classical lossless methods.
Binary MP4 files are similarly incompressible: zstd achieves only 1.26$\times$
on video content.

\subsection{Per-Layer Latency Breakdown}

\begin{table}[t]
\caption{Median per-layer latency [IQR] over D1 by file type (ms per file),
Pipeline~B (raw bytes $\to$ AES-256-GCM $\to$ ML-KEM-1024). VAE and I/O are not
separately isolated in this measurement.}
\label{tab:latency}
\centering\small
\begin{tabular}{lllll}
\toprule
File type & $L_\text{AES}$ (ms) & $L_\text{KEM}$ (ms) & $L_\text{total}$ (ms) & Crypto (\%) \\
\midrule
Images (JPEG) & 6.601 & 24.833 & 31.370 [$\pm$3.5] & 100.0 \\
Binary (MP4) & 40.270 & 29.632 & 71.645 [$\pm$79.6] & 97.6 \\
\midrule
All D1 & 7.126 & 26.496 & 34.235 [$\pm$31.0] & 98.2 \\
\bottomrule
\end{tabular}
\end{table}

Table~\ref{tab:latency} reports the measured per-file AES and ML-KEM latency for
Pipeline~B (no compression stage). Since only cryptographic operations are timed,
the \emph{Crypto~(\%)} column reflects the fraction of the \emph{measured} latency
attributable to crypto --- by construction nearly 100\%. In a complete end-to-end
pipeline that also measures VAE inference time ($\sim$6--70~ms per image depending
on model and hardware) and file I/O time, the crypto fraction is expected to fall
well below 50\%.

The large IQR for binary (MP4) files ($\pm$79.6~ms) reflects the high variance
in file sizes within the binary subset, ranging from small clips (few MB) to
large videos ($>$100~MB). The KEM latency is essentially constant per-file
(24.8--29.6~ms), as expected since it is independent of payload size.

\subsection{Compression and Round-Trip Fidelity}

\begin{table}[t]
\caption{Compression ratio and round-trip fidelity by file type.}
\label{tab:fidelity}
\centering\small
\begin{tabular}{p{1.6cm}p{0.9cm}p{1.2cm}p{1.2cm}p{1.6cm}}
\toprule
File type & CR & SSIM & PSNR (dB) & Notes \\
\midrule
Images (JPEG) & $\sim$6$\times$ (VAE) & 14.48~dB (PSNR) & 14.48 & 2-epoch prototype; lossy \\
Images (lossless) & $\sim$1.0$\times$ & 1.00 & 99.0 & JPEG already compressed \\
Binary (MP4) & $\sim$1.0--1.3$\times$ & N/A & N/A & High-entropy; minimal compression \\
\bottomrule
\end{tabular}
\end{table}

The VAE achieves 6$\times$ compression on the resized 32$\times$32 input but at
the cost of significant information loss (PSNR~$=~14.48$~dB after only 2 training
epochs). Binary MP4 files cannot be meaningfully compressed by either the VAE or
classical compressors, and are in fact slightly expanded (by 0--3\%) when
classical compression is naively applied. This is an acknowledged design
limitation; a file-type-aware bypass routing binary data directly to AES-GCM is
the correct engineering remedy and is designated as priority future work.

\subsection{Ablation Study}

\begin{table}[t]
\caption{Ablation study: measured throughput and security properties by pipeline
configuration. \emph{zstd} is used as the compression proxy for config A and
\emph{gzip} for config D (VAE end-to-end throughput is reported separately).}
\label{tab:ablation}
\centering\small
\begin{tabular}{p{2.2cm}rrll}
\toprule
Config & Img (MB/s) & Bin (MB/s) & PQ & Compression \\
\midrule
A: zstd+AES+KEM & 121.7 & 220.8 & Yes & zstd (proxy) \\
B: raw+AES+KEM & 203.8 & 578.3 & Yes & None \\
C: raw+AES & 993.6 & 912.8 & No & None \\
D: gzip+AES+KEM & 23.9 & 24.6 & Yes & gzip (slow) \\
\bottomrule
\end{tabular}
\end{table}

Four pipeline configurations are benchmarked (Table~\ref{tab:ablation}):

\begin{itemize}
  \item \textbf{Config~A} (zstd+AES+KEM): zstd compression achieves modest speed
  gains on binary data but the KEM overhead brings image throughput to 121.7~MB/s.
  \item \textbf{Config~B} (raw+AES+KEM): the baseline post-quantum pipeline
  without compression. Binary throughput reaches 578.3~MB/s.
  \item \textbf{Config~C} (raw+AES only): removing ML-KEM yields the highest
  throughput ($\sim$993~MB/s for images, $\sim$913~MB/s for binary), at the cost
  of losing post-quantum key security.
  \item \textbf{Config~D} (gzip+AES+KEM): gzip compression on already-compressed
  MP4/JPEG data collapses throughput to $\sim$24~MB/s due to gzip's high
  latency on incompressible content.
\end{itemize}

The comparison A vs.\ B isolates the cost of compression; B vs.\ C isolates the
ML-KEM overhead; C vs.\ A provides the net effect of adding both KEM and
compression. Config~D demonstrates the danger of applying an inappropriate
compressor to high-entropy data.

\subsection{Scalability}

\begin{table}[t]
\caption{Scalability results on D1 (3.79~GB, 126 files).}
\label{tab:scalability}
\centering\small
\begin{tabular}{rrrr}
\toprule
Size (GB) & Throughput (MB/s) & Latency (s) & Files \\
\midrule
0.60 & 819.4 & 0.75 & 53 \\
1.16 & 1056.3 & 0.95 & 71 \\
1.66 & 1014.6 & 1.27 & 78 \\
2.34 & 949.6 & 1.82 & 90 \\
2.87 & 834.9 & 2.46 & 108 \\
3.30 & 961.4 & 2.34 & 111 \\
3.79 & 954.2 & 2.58 & 126 \\
\bottomrule
\end{tabular}
\end{table}

Linear regression over the throughput-vs-size curve yields slope $\beta = -0.09$~MB/s/GB
and $R^2 < 0.01$. The near-zero $R^2$ indicates that the linear model explains
essentially none of the throughput variance at this scale: throughput fluctuates
between 820 and 1,056~MB/s with no significant trend across 0.6--3.79~GB.

This is a positive result: there is \emph{no evidence of throughput collapse} in
the measured range. The high variance ($\pm$74~MB/s around a mean of
941.5~MB/s) reflects file-type composition heterogeneity across cumulative batches
--- image-heavy batches process faster per byte than binary-heavy batches due to
the fixed KEM overhead per file.

\textbf{Comparison with original paper projection:} The original paper projected a
slope of $-2.87$~MB/s/GB and $R^2 = 0.824$ based on synthetic data modelling a
34.68~GB dataset. Our real measurements over 3.79~GB show a much flatter profile.
Evaluation on a larger dataset (>30~GB) is required to determine whether
throughput degradation becomes significant at scale; this is designated as future
work.

% ─────────────────────────────────────────────────────────────
\section{Discussion}

\subsection{Cryptographic Overhead is Modest but Not Negligible}

The measured ML-KEM-1024 overhead of $\sim$47~ms per key exchange is significant
for small files (a 7~KB JPEG image takes $\sim$7~ms for AES but $\sim$25~ms for
KEM decapsulation), but amortises favourably for large binary files ($>$10~MB),
where AES dominates. A session-key reuse strategy (encrypting multiple files under
a single ML-KEM key exchange) would reduce the per-file KEM cost to near-zero at
the cost of forward secrecy granularity.

\subsection{VAE Compression: Utility and Limitations}

The PSNR of 14.48~dB after 2 training epochs reflects a prototype implementation
only. Production training on 20~epochs with a diverse image corpus and perceptual
loss terms is expected to raise reconstruction quality significantly. Even then,
the VAE's 32$\times$32 input constraint means it is encoding a thumbnail, not the
full-resolution image. For applications requiring full-resolution reconstruction,
a modern learned codec (e.g., VQ-VAE-2, HiFiC) operating at full resolution
should replace the current architecture.

The binary data (MP4 files) cannot benefit from the image-trained VAE. Config~B
(raw+AES+KEM) achieves 578~MB/s on binary files, while Config~A (zstd+AES+KEM)
achieves only 221~MB/s. The compression stage provides no throughput benefit for
high-entropy binary content.

\subsection{Comparison with Related Work}

Our ablation study (Table~\ref{tab:ablation}) provides the first layer-isolated
quantification of throughput for a combined VAE + AES-GCM + ML-KEM configuration.
The comparison A vs.\ B shows that adding ML-KEM reduces binary throughput from
578.3 to 220.8~MB/s --- a 2.6$\times$ slowdown. This is substantially more than
the overhead reported by Fu et al.\ \cite{fu} for classical RSA-2048 key exchange
(approximately 1.5$\times$ overhead), reflecting the higher pure-Python ML-KEM
cost vs.\ an optimised RSA implementation. Native C-compiled liboqs on Linux is
expected to reduce this gap.

\subsection{Deployment Taxonomy}

We propose the following deployment taxonomy:
\begin{itemize}
  \item \textbf{Appropriate:} Image similarity search, content-based media
  retrieval, streaming preview generation, privacy-preserving approximate analytics.
  \item \textbf{Inappropriate:} Archival storage, financial transaction records,
  medical imaging at diagnostic quality, any application governed by data integrity
  regulations requiring bit-exact reconstruction.
\end{itemize}

For inappropriate use cases, routing binary and tabular data directly to
AES-256-GCM (Config~B) provides full throughput without VAE overhead.

\subsection{Limitations}

\begin{enumerate}
  \item \textbf{CPU-only, pure-Python implementation:} Absolute latency figures
  are not directly comparable to production bare-metal deployments with
  native-compiled liboqs.
  \item \textbf{VAE prototype:} 2-epoch training is insufficient for production
  use. The compression layer requires significant further development.
  \item \textbf{No large-scale (D2) evaluation:} Scalability beyond 3.79~GB was
  not evaluated in this study. Throughput behaviour at $>$30~GB remains
  uncharacterised.
  \item \textbf{Informal security composition:} The KEM-DEM composition argument
  is not a formal proof in the random oracle model.
  \item \textbf{No adversarial robustness evaluation:} Robustness of the VAE
  decoder to adversarial perturbations is not evaluated.
\end{enumerate}

% ─────────────────────────────────────────────────────────────
\section{Conclusion}

We have presented, implemented, and evaluated a reproducible hybrid security
pipeline integrating VAE compression, AES-256-GCM authenticated encryption, and
ML-KEM/Kyber-1024 post-quantum key encapsulation on 3.79~GB of real heterogeneous
data (126 files: 75 JPEG images, 50 MP4 binaries, 1 JSON text file). The key
findings are:

\begin{enumerate}
  \item \textbf{NIST CAVP correctness:} 775/775 test vectors passed (100\%) for
  both AES-256-GCM and ML-KEM-1024.
  \item \textbf{ML-KEM overhead is fixed per-file:} 47.4~ms per key exchange
  (pure-Python), modest for large files, significant for small files.
  \item \textbf{Classical compression fails on high-entropy data:} MP4/JPEG files
  are already near-incompressible; gzip collapses throughput to 24~MB/s.
  \item \textbf{No throughput collapse at measured scale:} Throughput fluctuates
  around 941~MB/s with no downward trend from 0.6 to 3.79~GB ($\beta = -0.09$
  MB/s/GB, $R^2 < 0.01$).
  \item \textbf{VAE requires further development:} 14.48~dB PSNR after 2 epochs
  is prototype-level; 20-epoch training with perceptual loss is the immediate
  next step.
\end{enumerate}

Future work will prioritise: (i)~bare-metal Linux evaluation with native liboqs
to remove pure-Python overhead; (ii)~a formal KEM-DEM composition proof; (iii)~an
entropy-aware file-type router providing a lossless bypass for binary data;
(iv)~VAE architecture improvements (perceptual loss, adversarial training, full
resolution) to raise PSNR above 25~dB; (v)~large-scale ($>$30~GB) scalability
evaluation; and (vi)~extension to continuous data streams (video, IoT telemetry)
as a step toward an intelligent adaptive security pipeline.

\section*{Declarations}

\textbf{Conflict of interest:} The authors declare no conflict of interest.

\textbf{Ethics statement:} This research uses only real files provided by the
authors. No personal data of third parties is processed.

\textbf{Data and code availability:} The pipeline source code and CSV result
files are available at [repository URL]. Full reproduction instructions are
provided in the repository README.

\textbf{AI assistance:} Large language model assistance (Claude, Anthropic) was
used for English grammar editing, structural revision, and code development
support. All scientific content, experimental design, results, and conclusions are
the sole responsibility of the authors.

% ─────────────────────────────────────────────────────────────
\begin{thebibliography}{99}

\bibitem{natacad2019}
National Academies of Sciences, Engineering, and Medicine,
\textit{Quantum Computing: Progress and Prospects}.
Washington, DC: The National Academies Press, 2019.
doi: 10.17226/25196.

\bibitem{bernstein2017}
D.~J. Bernstein and T.~Lange,
``Post-quantum cryptography,''
\textit{Nature}, vol.~549, no.~7671, pp.~188--194, Sep.~2017.
doi: 10.1038/nature23461.

\bibitem{laney2001}
D.~Laney,
``3D Data Management: Controlling Data Volume, Velocity, and Variety,''
META Group Research Note, Feb.~2001.

\bibitem{gandomi2015}
A.~Gandomi and M.~Haider,
``Beyond the hype: Big data concepts, methods, and analytics,''
\textit{International Journal of Information Management},
vol.~35, no.~2, pp.~137--144, Apr.~2015.
doi: 10.1016/j.ijinfomgt.2014.10.007.

\bibitem{grover1996}
L.~K. Grover,
``A fast quantum mechanical algorithm for database search,''
in \textit{Proc. 28th Annual ACM Symposium on Theory of Computing (STOC)},
Philadelphia, PA, 1996, pp.~212--219.
doi: 10.1145/237814.237866.

\bibitem{dworkin2007}
M.~Dworkin,
``NIST Special Publication 800-38D: Recommendation for Block Cipher Modes
of Operation: Galois/Counter Mode (GCM) and GMAC,''
National Institute of Standards and Technology, Nov.~2007.
doi: 10.6028/NIST.SP.800-38D.

\bibitem{aviram2016}
N.~J. Aviram et al.,
``DROWN: Breaking TLS using SSLv2,''
in \textit{Proc. 25th USENIX Security Symposium},
Austin, TX, 2016, pp.~689--706.

\bibitem{fips203}
National Institute of Standards and Technology,
``FIPS 203: Module-Lattice-Based Key-Encapsulation Mechanism Standard,''
Federal Information Processing Standards Publication 203, Aug.~2024.
doi: 10.6028/NIST.FIPS.203.

\bibitem{zhang}
[Zhang et al.\ reference --- please supply full citation details from Zotero.]

\bibitem{fu}
[Fu et al.\ reference --- please supply full citation details from Zotero.]

\bibitem{nugroho}
[Nugroho et al.\ reference --- please supply full citation details from Zotero.]

\bibitem{thenmozhi}
[Thenmozhi et al.\ reference --- please supply full citation details from Zotero.]

\bibitem{alsubai}
[Alsubai et al.\ reference --- please supply full citation details from Zotero.]

\bibitem{wanglo}
B.~Wang and K.-T. Lo,
``Autoencoder-based joint image compression and encryption,''
\textit{Journal of Information Security and Applications},
vol.~80, p.~103680, Feb.~2024.
doi: 10.1016/j.jisa.2023.103680.

\bibitem{fall}
[Fall et al.\ --- verify publication status. If still a preprint:
J.~Fall, ``SoK: Systematizing Hybrid Strategies for Post-Quantum Key Exchange,''
IACR ePrint 2025/2052, preprint, 2025.]

\bibitem{shoup2001}
V.~Shoup,
``A Proposal for an ISO Standard for Public Key Encryption,''
IACR ePrint 2001/112, 2001.
Available: \url{https://eprint.iacr.org/2001/112}.

\bibitem{cramer2003}
R.~Cramer and V.~Shoup,
``Design and Analysis of Practical Public-Key Encryption Schemes Secure against
Adaptive Chosen Ciphertext Attack,''
\textit{SIAM Journal on Computing},
vol.~33, no.~1, pp.~167--226, Jan.~2003.
doi: 10.1137/S0097539702403773.

\end{thebibliography}

\end{document}
